% ---
% titre: 9P - Nombres naturels et décimaux
% theme: NO
% chapitre: Nombres naturels et décimaux
% sous_chapitre:
% niveau: [9e]
% type: PLANIF
% tags: []
% source: latex
% ---

\header{Nombres naturels et décimaux - Objectifs}

\def\arraystretch{1.2}
\begin{tabularx}{\linewidth}{|>{\centering\arraybackslash}p{3cm}|>{\centering\arraybackslash}p{3cm}|X|}
\hline
\cellcolor{themeNO!30} Livre & \cellcolor{themeNO!30} Fiches & \cellcolor{themeNO!30} \hfill Aide-mémoire\hfill\null\\
\hline
&&Généralité: pp. 24-26\\
pp. 23-43&pp. 1-38&Nombres naturels: pp. 26-30\\
&&Nombres décimaux: pp. 35-41\\
\hline
\end{tabularx}
\def\arraystretch{1}

\def\arraystretch{1.6}
\begin{tabularx}{\textwidth}{|>{\raggedright\arraybackslash}X|>{\raggedright\arraybackslash}X|}
\hline
\cellcolor{themeNO!30}Compétences visées & \cellcolor{themeNO!30}Utilité dans la vie de tous les jours \\
\hline
Mobiliser les savoirs, savoir-faire et savoir être pour résoudre des problèmes en lien avec les nombres naturels et décimaux.
&\vspace{-1em}
\begin{itemize}[leftmargin=*, itemsep=2pt, topsep=0pt, partopsep=0pt, parsep=0pt, label=$\triangleright$]
    \item Évaluer la température
    \item Se repérer entre les étages d'un bâtiment
    \item Évaluer les prix, les dépenses, le salaire
\end{itemize} \\
\hline

Développer des stratégies de recherche afin de résoudre des problèmes.
&\vspace{-1em}
\begin{itemize}[leftmargin=*, itemsep=2pt, topsep=0pt, label=$\triangleright$]
    \item Découper un projet en sous-tâches
    \item Choisir l'ordre des tâches à effectuer
\end{itemize} \\
\hline

Présenter sa démarche de façon claire (explications, grandeur cherchée, unités, calculs, réponse).
&\vspace{-1em}
\begin{itemize}[leftmargin=*, itemsep=2pt, topsep=0pt, label=$\triangleright$]
    \item Collaborer avec d'autres personnes
    \item Expliciter sa compréhension d'un sujet
    \item Argumenter
\end{itemize} \\
\hline

Développer des stratégies d'apprentissage pour s'approprier puis mobiliser les savoirs et savoir-faire.
&\vspace{-1em}
\begin{itemize}[leftmargin=*, itemsep=2pt, topsep=0pt, label=$\triangleright$]
    \item Apprendre une nouvelle activité %(art, sport, \ldots)
    \item S'adapter à un nouvel environnement %(travail, voyage, \ldots)
    \item Se former en autonomie %(tutoriel, lecture, \ldots)
\end{itemize} \\
\hline

Vérifier sa solution ou au moins sa plausibilité.
&\vspace{-1em}
\begin{itemize}[leftmargin=*, itemsep=2pt, topsep=0pt, label=$\triangleright$]
    \item S'autocritiquer
    \item Évaluer la qualité de son travail
\end{itemize} \\
\hline
\end{tabularx}
\def\arraystretch{1}

\newpage
\null
\vfill
%---------------
\def\arraystretch{1.6}
\begin{tabularx}{\textwidth}{|>{\raggedright\arraybackslash}p{5.2cm}|>{\raggedright\arraybackslash}X|>{\centering\arraybackslash}p{0.8cm}|>{\centering\arraybackslash}p{0.8cm}|>{\centering\arraybackslash}p{0.8cm}|}
\hline
\rowcolor{themeNO!30}\multicolumn{2}{|l|}{Objectifs} &\multicolumn{3}{c|}{Auto-évaluation} \\
\rowcolor{themeNO!30}\multicolumn{2}{|l|}{} & \Sadey[1.5]& \Neutrey[1.5] & \Smiley[1.5]  \\\hline
 
\multirow{6}{=}{Reconnaître et utiliser certaines propriétés des nombres naturels:
\begin{itemize}
    \item critères de divisibilité, multiples et diviseurs communs
    \item ppmc, pgdc, nombres premiers, produit de facteurs
\end{itemize}}
& Maîtriser les livrets jusqu'à 12 & & & \\
\cline{2-5}
& Déterminer si un nombre est premier & & & \\
\cline{2-5}
& Connaître les nombres premiers jusqu'à 23 & & & \\
\cline{2-5}
& Déterminer si un nombre est divisible par 2, 3, 4, 5, 6, 9, 10, 25, 50 et 100 à l'aide des critères de divisibilité & & & \\
\cline{2-5}
& Décomposer un nombre en produit de facteurs premiers& & & \\
\cline{2-5}
& Trouver le ppmc de plusieurs nombres:
\begin{itemize}[label=$\triangleright$]
    \item en établissant la liste de leurs multiples
    \item en utilisant la décomposition en produit de facteurs premiers
\end{itemize} & & & \\
\cline{2-5}
 & Trouver le pgdc de plusieurs nombres:
\begin{itemize}[label=$\triangleright$]
    \item en établissant la liste de leurs diviseurs
    \item en utilisant la décomposition en produit de facteurs premiers
\end{itemize} & & & \\
\hline
 
\multirow{2}{=}{Connaître et utiliser différentes écritures d'un même nombre (y compris sous forme de puissances)}
& Connaître et savoir utiliser le vocabulaire en lien avec les nombres et les opérations && & \\
\cline{2-5}
& Écrire et lire des nombres écrits sous forme de puissance (carré ou cube) & & & \\
\hline
 
\multirow{2}{=}{Connaître et utiliser les priorités des opérations (y compris parenthèses)}
& Effectuer des calculs en respectant les priorités des opérations. & & & \\
\cline{2-5}
& Comprendre le rôle des parenthèses et les utiliser à bon escient dans un calcul& & & \\
\hline
 
Utiliser les propriétés des opérations (addition, soustraction, multiplication, division) pour calculer efficacement et estimer
& Mobiliser les propriétés des opérations (commutativité, associativité, distributivité) pour:
\begin{itemize}[label=$\triangleright$]
    \item calculer mentalement, ex.~: $57-38$ \textbar\ $15 \cdot 40$
    \item organiser un calcul pour le simplifier, ex.~: $47 \cdot 8 - 7 \cdot 8$ \textbar\ $4 \cdot 3{,}2 \cdot 25$
\end{itemize} & & & \\
\hline
 \end{tabularx}
 \def\arraystretch{1}
 
%------------------------------------------------
\newpage
\header{Nombres naturels et décimaux - Feuille de route}

\newcommand{\rref}[1]{\textcolor{themeColor}{#1}}
 
% Case à cocher
\newcommand{\case}{\raisebox{-0.1ex}{\Large$\square$}}
 
% Marqueurs obligatoire / facultatif (forme, pas seulement couleur)
\newcommand{\oblig}{\tikz[baseline=-0.5ex]{\fill[themeColor] (0,0) circle (0.09);}}
\newcommand{\faclt}{\tikz[baseline=-0.5ex]{\draw[gray!70,line width=0.6pt] (0,0) circle (0.09);}}
 
 \vfill
\begin{tikzpicture}[
circ/.style={circle, fill=themeColor, text=white, font=\bfseries\Large, minimum size=11mm, inner sep=0pt},
  box/.style={rectangle, draw=themeColor!35!white, rounded corners=2pt, fill=themeColor!8!white, inner sep=3mm, text width=14.7cm, align=left, anchor=north west}]
 
% ---------- Étape 1 ----------
\node[circ, anchor=north] (c1) at (0,0) {1};
\node[box] (b1) at ($(c1.north)+(1.2cm,0)$) {%
\textbf{\large Critères de divisibilité}\\[4pt]
{\small\begin{tabularx}{10cm}{cXp{3cm}p{1cm}}
\case & \rref{NO 3}       & Fiche, p.\,3 & \oblig \\
\case & \rref{Problème 3} & VLM, p.\,4    & \oblig \\
\case & \rref{NO 28}      & Livre, p.\,29 & \oblig \\
\end{tabularx}}
};
 
% ---------- Étape 2 ----------
\coordinate (y2) at ($(b1.south west)+(0,-3mm)$);
\node[circ, anchor=north] (c2) at (c1 |- y2) {2};
\node[box] (b2) at (y2) {%
\textbf{\large Nombres premiers et décomposition}\\[4pt]
{\small\begin{tabularx}{10cm}{cXp{3cm}p{1cm}}
\case & \rref{NO 50}         & Livre, p.\,33 & \oblig \\
\case & \rref{Problème 5 A.} & VLM, p.\,5    & \oblig \\
\case & Problème 5 B.        & VLM, p.\,5    & \faclt \\
\case & Problème 5 C.        & VLM, p.\,5    & \faclt \\
\end{tabularx}}
};
 
% ---------- Étape 3 ----------
\coordinate (y3) at ($(b2.south west)+(0,-3mm)$);
\node[circ, anchor=north] (c3) at (c1 |- y3) {3};
\node[box] (b3) at (y3) {%
\textbf{\large PPMC et PGDC}\\[4pt]
{\small\begin{tabularx}{10cm}{cXp{3cm}p{1cm}}
\case & \rref{NO 52} & Livre, p.\,33 & \oblig \\
\case & \rref{NO 53} & Livre, p.\,33 & \oblig \\
\case & \rref{NO 54} & Livre, p.\,33 & \oblig \\
\case & \rref{Problème 10 C., H., J.} & VLM, p.\,7-8 & \oblig \\
\case & Problème 10 Autres           & VLM, p.\,7-8 & \faclt \\
\end{tabularx}}
};
 
% ---------- Étape 4 ----------
\coordinate (y4) at ($(b3.south west)+(0,-3mm)$);
\node[circ, anchor=north] (c4) at (c1 |- y4) {4};
\node[box] (b4) at (y4) {%
\textbf{\large Puissances}\\[4pt]
{\small\begin{tabularx}{10cm}{cXp{3cm}p{1cm}}
\case & \rref{NO 62} & Fiche, p.\,21 & \oblig \\
\case & \rref{NO 63} & Fiche, p.\,21 & \oblig \\
\case & \rref{NO 64} & Fiche, p.\,22 & \oblig \\
\end{tabularx}}
};
 
% ---------- Étape 5 ----------
\coordinate (y5) at ($(b4.south west)+(0,-3mm)$);
\node[circ, anchor=north] (c5) at (c1 |- y5) {5};
\node[box] (b5) at (y5) {%
\textbf{\large Priorité des opérations}\\[4pt]
{\small\begin{tabularx}{10cm}{cXp{3cm}p{1cm}}
\case & \rref{NO 78}  & Fiche, p.\,26 & \oblig \\
\case & \rref{NO 82}  & Livre, p.\,40 & \oblig \\
\case & \rref{NO 100} & Fiche, p.\,37 & \oblig \\
\case & \rref{NO 89}  & Livre, p.\,43 & \oblig \\
\case & NO 98 & Fiche, p.\,36    & \faclt \\
\case & NO 99 & Fiche, p.\,37    & \faclt \\
\case & NO 75 & Fiche, p.\,24    & \faclt \\
\case & NO 76 & Fiche, p.\,24-25 & \faclt \\
\end{tabularx}}
};
 
% ---------- Étape 6 ----------
\coordinate (y6) at ($(b5.south west)+(0,-3mm)$);
\node[circ, anchor=north] (c6) at (c1 |- y6) {6};
\node[box] (b6) at (y6) {%
\textbf{\large Exercices \guillemotleft\,Défi\,\guillemotright}\\[4pt]
{\small\begin{tabularx}{10cm}{cXp{3cm}p{1cm}}
\case & NO 43      & Livre, p.\,32 & \faclt \\
\case & NO 73      & Livre, p.\,35 & \faclt \\
\case & NO 60      & Livre, p.\,38 & \faclt \\
\case & NO 15      & Fiche, p.\,4  & \faclt \\
\case & NO 51      & Livre, p.\,33 & \faclt \\
\case & Problème 2 & VLM, p.\,2    & \faclt \\
\end{tabularx}}
};
 
% ---------- Ligne de parcours (derrière les cercles) ----------
\begin{pgfonlayer}{background}
\draw[themeColor, line width=2pt] (c1) -- (c2) -- (c3) -- (c4) -- (c5) -- (c6);
\end{pgfonlayer}
 
\end{tikzpicture}
 
 \vfill
\begin{center}
\small \case~à faire \qquad \oblig~\rref{obligatoire} \qquad \faclt~facultatif (consolidation / pour aller plus loin)
\end{center}
\vspace{1em}

%------------------------------------------------
\newtcolorbox{regle}[1]{enhanced, sharp corners, boxrule=0pt, interior hidden, frame hidden,borderline west={2.2pt}{0pt}{themeColor},left=3mm, right=2mm, top=2mm, bottom=2mm,title={\textbf{\textcolor{themeColor}{#1}}}, fonttitle=\bfseries,attach title to upper={\\},}
\newtcolorbox{astuce}[1]{enhanced, sharp corners, halign=justify, valign=top, colback=white, colframe=themeColor, boxrule=2.2pt, left=3mm, right=2mm, top=2mm, bottom=2mm,title={\textbf{\textcolor{themeColor}{#1}}}, fonttitle=\bfseries\small,attach title to upper={\\},}

\newpage
\header{Nombres naturels et décimaux - Théorie}

\subsection*{1. Multiplication et puissance}

\begin{regle}{Définition}
Une \textcolor{themeNO}{multiplication} est une répétition d'additions. Elle se note avec un point médiant $\cdot$.
\end{regle}

\begin{center}$\underset{5\ fois}{\underbrace{3+3+3+3+3}}=5\cdot 3$\end{center}

\begin{regle}{Définition}
Une \textcolor{themeNO}{puissance} est une répétition de multiplications. Elle se note avec un exposant.
\end{regle}

\begin{center}$\underset{5\ fois}{\underbrace{3\cdot 3\cdot 3\cdot 3\cdot3}}=3^{5}$\end{center}

\begin{astuce}{Vocabulaire}
$3^5$ se lit \og 3 \textcolor{themeNO}{à la puissance} 5 \fg.\\[2pt]
Lorsqu'un nombre est à la puissance 2, on dit qu'il est au \textcolor{themeNO}{carré}.\\[2pt]
Lorsqu'un nombre est à la puissance 3, on dit qu'il est au \textcolor{themeNO}{cube}.
\end{astuce}

\vfill
\subsection*{2. Priorité des opérations}

Comme les multiplications sont des additions cachées et que les puissances sont des multiplications cachées, les opérations doivent être effectuées dans un certain ordre pour être correctes.

\begin{center}
$7+5\cdot 2=7+\underset{2\ fois}{\underbrace{5+5}}=17$
\end{center}

La priorité des opérations respecte le \textcolor{themeNO}{PEMDAS}.

\def\arraystretch{1.6}
\begin{tabularx}{\linewidth}{|p{2cm}|>{\centering\arraybackslash}p{2cm}|X|}
\hline
1ère & \textcolor{themeNO}{P} & Parenthèses\\\hline
2ème & \textcolor{themeNO}{E} & Exposants (puissances)\\\hline
3ème & \textcolor{themeNO}{MD} & Multiplications et divisions\\\hline
4ème & \textcolor{themeNO}{AS} & Additions et soustractions\\\hline
\end{tabularx}
\def\arraystretch{1}

\newpage
Exemple :

\vspace{-10pt}\begin{center}$3\cdot 50+\left(10-8\right)^3-16\div4$\end{center}

\vspace{-30pt}\begin{align*}
3\cdot 50+\left(\textcolor{themeNO}{10-8}\right)^3-16\div4&=&&\text{On résout le contenu des parenthèses (\textcolor{themeNO}{P}).}\\
3\cdot 50+\textcolor{themeNO}{2^3}-16\div4&=   &&\text{On effectue les exposants (\textcolor{themeNO}{E}).}\\
\textcolor{themeNO}{3\cdot 50}+8-\textcolor{themeNO}{16\div4}&=&&\text{On effectue multiplications et divisions (\textcolor{themeNO}{MD}) de gauche à droite.}\\
\textcolor{themeNO}{150+8-4}&=&&\text{On effectue additions et soustractions (\textcolor{themeNO}{AS}) de gauche à droite.}\\
\textcolor{themeNO}{154}&&&
\end{align*}

\vfill
\subsection*{3. Multiples et diviseurs}

\begin{regle}{Définition}
Un nombre entier $N$ est un \textcolor{themeNO}{multiple} d'un autre nombre entier $n_1$ s'il existe un troisième nombre entier $n_2$ tel que $n_1\cdot n_2=N$.
\end{regle}

\begin{center}15 est un multiple de 5, car il existe un nombre (3) tel que $5\cdot 3=\underset{multiple}{\underbrace{15}}$\end{center}

\begin{regle}{Définition}
Un nombre entier $n$ est un \textcolor{themeNO}{diviseur} d'un autre nombre entier $N$ si la division de $N$ par $n$ donne un nombre entier. 
\end{regle}

\begin{center}12 est un diviseur de 72, car $72\div12=\underset{entier}{\underbrace{6}}$\end{center}

\vfill
\subsection*{4. Critères de divisibilité }

Il n'est pas toujours facile de déduire si un nombre est divisible par un autre. Voici les critères permettant de savoir si un nombre est divisible par 2, 3, 4, 5, 6, 8, 9, 10, 25, 50 ou 100.

\vspace{10pt}
\def\arraystretch{1.1}
\begin{tabularx}{\linewidth}{|c|X|}
\hline
Divisible par?&Critère\\\hline
2&Le nombre est pair, il finit donc par 0,2,4,6,8).\\\hline
3&La somme des chiffres formant le nombre est divisible par 3.\\\hline
4&Les deux derniers chiffres du nombre sont divisibles par 4.\\\hline
5&Le nombre finit par 0 ou 5.\\\hline
6&Le nombre est simultanément divisible par 2 et 3.\\\hline
9&La somme des chiffres formant le nombre est divisible par 3.\\\hline
10&Le nombre finit par 0.\\\hline
25&Le nombre finit par 00, 25, 50 ou 75.\\\hline
50&Le nombre finit par 00 ou 50.\\\hline
100&Le nombre finit par 00.\\\hline
\end{tabularx}
\def\arraystretch{1}

\newpage
\subsection*{5. Nombres Premiers}

\begin{regle}{Définition}
Un nombre premier est un nombre qui n'est divisible que par 1 et par lui-même. Il en existe une infinité.
\end{regle}

\begin{center}
23 n'est divisible que par 1 et par 23. C'est donc un nombre premier.\\
49 est divisible par 7. Ce n'est donc pas un nombre premier.
\end{center}

\vspace{20pt}\begin{astuce}{Nombres premiers à connaître par coeur}
\hfill 2 \hfill 3 \hfill 5	 \hfill 7\hfill 11\hfill 13\hfill 17\hfill 19\hfill 23
\end{astuce}

\vfill
\subsection*{6. Décomposition en facteurs premiers}

N'importe quel nombre entier peut être écrit comme une multiplication de nombres premiers. Cette multiplication est unique pour chaque nombre.

\begin{center}
$45=3\cdot 3\cdot 5=3^2\cdot 5$\\
\end{center}

Pour trouver la décomposition en facteurs premiers d'un nombre, il faut le diviser autant de fois que possible par les nombres premiers dans l'ordre.

\vspace{30pt}
Exemples: 

\vspace{-5pt}\def\arraystretch{1.7}
\begin{tabularx}{0.3\linewidth}{R|X}
 \multicolumn{2}{c}{Décomposition de 2310}\\
2310&2\\
1155&3\\
385&5\\
77&7\\
11&11\\
1&\\
\multicolumn{2}{c}{$2310=2\cdot3\cdot5\cdot7\cdot11$}\\
\end{tabularx}
\hfill
\begin{tabularx}{0.3\linewidth}{R|X}
 \multicolumn{2}{c}{Décomposition de 264}\\
264&2\\
132&2\\
66&2\\
33&3\\
11&11\\
1&\\
\multicolumn{2}{c}{$264=2^3\cdot3\cdot11$}\\
\end{tabularx}
\hfill
\begin{tabularx}{0.3\linewidth}{R|X}
 \multicolumn{2}{c}{Décomposition de 350}\\
350&2\\
175&5\\
35&5\\
7&7\\
1&\\
&\\
\multicolumn{2}{c}{$2\cdot5^2\cdot7$}\\
\end{tabularx}
\def\arraystretch{1}

\newpage
\subsection*{7. Plus petit multiple commun (PPMC)}

\begin{regle}{Définition}
Le \textcolor{themeNO}{plus petit multiple commun} de $n$ entiers non nuls est le plus petit entier strictement positif multiple simultanément de ces $n$ entiers.
\end{regle}

\begin{center}
ppmc(45,54) =270
\end{center}

\vspace{40pt}
Pour déterminer le ppmc de deux nombres, il existe deux méthodes.

\vspace{20pt}
La \textcolor{themeNO}{méthode de la liste} consiste à lister les multiples de chacun des nombres jusqu'à trouver un nombre qui apparaît dans chacune des listes.

$M\left(45\right)=\left\{45, 90, 135, 180, 225, 270, 315, 360, 405, \ldots \right\}$\\
$M\left(54\right)=\left\{54, 108, 162, 216, 270, 324, 378, 432, \ldots \right\}$\\
ppmc$(45,54) =270$


\vspace{20pt}
La \textcolor{themeNO}{méthode de la décomposition en facteurs premiers} consiste à déterminer la décomposition en facteurs premiers de chaque nombre et pour chaque facteur, prendre sa plus grande puissance.

$45=3^{2}\cdot 5$ \quad\quad et \quad\quad $54=2\cdot 3^3$

Les facteurs présents sont 2, 3 et 5. Pour chaque facteur, nous prenons la plus grande puissance.

$2\cdot 3^3\cdot 5=270$\\
ppmc$(45,54) =270$

\newpage
\subsection*{8.  Plus grand diviseur commun (PGDC)}

\begin{regle}{Définition}
Le \textcolor{themeNO}{plus grand diviseur commun} de deux nombres entiers non nuls est le plus grand nombre entier qui divise ces deux nombres en même temps sans laisser de reste.
\end{regle}

\begin{center}
pgdc(24,60) =12
\end{center}

\vspace{40pt}
Pour déterminer le pgdc de deux nombres, il existe deux méthodes.

\vspace{20pt}
La \textcolor{themeNO}{méthode de la liste} consiste à lister les diviseurs de chacun des nombres et choisir le plus grand nombre qui apparaît dans les deux listes.

$D\left(24\right)=\left\{1, 2, 3, 4, 6, 8, 12, 24\right\}$\\
$D\left(60\right)=\left\{1, 2, 3, 4, 5, 6, 10, 12, 15, 20, 30, 60\right\}$\\
pgdc$(24,60) =12$

\vspace{20pt}
La \textcolor{themeNO}{méthode de la décomposition en facteurs premiers} consiste à déterminer la décomposition en facteurs premiers de chaque nombre et pour chaque facteur présent dans les deux décompositions, prendre sa plus petite puissance si elle existe.

$24=2^3\cdot 3$ \quad\quad et \quad\quad $60=2^{2}\cdot 3\cdot 5$

Les seuls facteurs présents dans les deux décompositions sont 2 et 3. Pour chaque facteur présent dans les deux décompositions, nous prenons la plus petite puissance.

$2^2\cdot 3=12$\\
pgdc$(24,60) =12$


%------------------------------------------------
\newpage
\header{Nombres naturels et décimaux - Lexique}

\newtcolorbox{terme}{%
  enhanced, sharp corners, boxrule=0pt, colback=themeColor!6!white,
  colframe=themeColor!6!white,
  borderline west={2.2pt}{0pt}{themeColor},
  left=3mm, right=2mm, top=2mm, bottom=2mm,
}
 
 \vfill
\begin{tcbraster}[raster columns=2, raster equal height, raster valign=top,raster column skip=2mm, raster row skip=2mm]
\begin{terme}
\textbf{\large Addition}\\[2pt]
{\small Opération qui consiste à réunir deux nombres (les termes) pour obtenir leur somme.}
\end{terme}
\begin{terme}
\textbf{\large Arrondi}\\[2pt]
{\small Valeur approchée d'un nombre, obtenue en supprimant les derniers chiffres selon une règle donnée.}
\end{terme}
\begin{terme}
\textbf{\large Associativité}\\[2pt]
{\small Propriété permettant de regrouper les termes ou facteurs d'une opération sans changer le résultat.}
\end{terme}
\begin{terme}
\textbf{\large Chiffre}\\[2pt]
{\small Symbole utilisé pour écrire un nombre (0 à 9 dans le système décimal).}
\end{terme}
\begin{terme}
\textbf{\large Commutativité}\\[2pt]
{\small Propriété permettant d'échanger l'ordre des termes ou facteurs d'une opération sans changer le résultat.}
\end{terme}
\begin{terme}
\textbf{\large Différence}\\[2pt]
{\small Résultat d'une soustraction.}
\end{terme}
\begin{terme}
\textbf{\large Distributivité}\\[2pt]
{\small Propriété qui permet de multiplier une somme ou une différence en multipliant chaque terme séparément.}
\end{terme}
\begin{terme}
\textbf{\large Dividende}\\[2pt]
{\small Nombre que l'on divise dans une division.}
\end{terme}
\begin{terme}
\textbf{\large Diviseur}\\[2pt]
{\small Nombre par lequel on divise ; nombre qui divise un autre nombre sans reste.}
\end{terme}
\begin{terme}
\textbf{\large Diviseur commun}\\[2pt]
{\small Nombre qui divise à la fois plusieurs nombres donnés.}
\end{terme}
\begin{terme}
\textbf{\large Division}\\[2pt]
{\small Opération qui permet de partager un nombre (le dividende) en parts égales.}
\end{terme}
\begin{terme}
\textbf{\large Élément neutre}\\[2pt]
{\small Nombre qui, utilisé dans une opération, ne change pas le résultat (0 pour l'addition, 1 pour la multiplication).}
\end{terme}
\begin{terme}
\textbf{\large Exposant}\\[2pt]
{\small Nombre placé en haut à droite d'un autre indiquant combien de fois celui-ci est multiplié par lui-même.}
\end{terme}
\begin{terme}
\textbf{\large Facteur}\\[2pt]
{\small Chacun des nombres que l'on multiplie entre eux.}
\end{terme}
\begin{terme}
\textbf{\large Nombre}\\[2pt]
{\small Notion mathématique servant à compter, mesurer ou repérer.}
\end{terme}
\begin{terme}
\textbf{\large Nombre naturel}\\[2pt]
{\small Nombre entier positif ou nul (0, 1, 2, 3, \ldots).}
\end{terme}
\begin{terme}
\textbf{\large Nombre décimal}\\[2pt]
{\small Nombre qui peut s'écrire avec une virgule (partie entière et partie décimale).}
\end{terme}
\begin{terme}
\textbf{\large Nombre premier}\\[2pt]
{\small Nombre naturel supérieur à 1 qui n'a que deux diviseurs~: 1 et lui-même.}
\end{terme}
\begin{terme}
\textbf{\large Multiple}\\[2pt]
{\small Résultat de la multiplication d'un nombre par un nombre naturel.}
\end{terme}
\begin{terme}
\textbf{\large Multiple commun}\\[2pt]
{\small Nombre qui est multiple de plusieurs nombres donnés à la fois.}
\end{terme}
\begin{terme}
\textbf{\large Multiplication}\\[2pt]
{\small Opération qui permet de calculer le produit de deux ou plusieurs facteurs.}
\end{terme}
\begin{terme}
\textbf{\large Ordre décroissant}\\[2pt]
{\small Rangement des nombres du plus grand au plus petit.}
\end{terme}
\begin{terme}
\textbf{\large Ordre croissant}\\[2pt]
{\small Rangement des nombres du plus petit au plus grand.}
\end{terme}
\begin{terme}
\textbf{\large Pgdc}\\[2pt]
{\small Plus grand diviseur commun à plusieurs nombres.}
\end{terme}
\begin{terme}
\textbf{\large Ppmc}\\[2pt]
{\small Plus petit multiple commun (non nul) à plusieurs nombres.}
\end{terme}
\begin{terme}
\textbf{\large Produit}\\[2pt]
{\small Résultat d'une multiplication.}
\end{terme}
\begin{terme}
\textbf{\large Puissance}\\[2pt]
{\small Écriture qui représente la multiplication répétée d'un même facteur.}
\end{terme}
\begin{terme}
\textbf{\large Quotient}\\[2pt]
{\small Résultat d'une division.}
\end{terme}
\begin{terme}
\textbf{\large Somme}\\[2pt]
{\small Résultat d'une addition.}
\end{terme}
\begin{terme}
\textbf{\large Soustraction}\\[2pt]
{\small Opération qui permet de calculer la différence entre deux nombres.}
\end{terme}
\begin{terme}
\textbf{\large Terme}\\[2pt]
{\small Chacun des nombres que l'on additionne.}
\end{terme}
\end{tcbraster}
